% Date: May 04, 2026
% Status: Resubmission - Cover Letter
%==============================================================================

\documentclass[11pt,a4paper]{letter}
\usepackage[utf8]{inputenc}
\usepackage[T1]{fontenc}
\usepackage{amsmath,siunitx}

\address{
Department of Physics\\
Faculty of Science\\
University of Douala, Cameroon\\
Email: steve.teguia@facsciences-uy1.cm
}

\signature{Steve Cabrel Teguia Kouam}

\begin{document}

\begin{letter}{
\textbf{Prof. Gregory Scholes}\\
Editor-in-Chief\\
The Journal of Physical Chemistry Letters\\
American Chemical Society
}

\opening{Dear Prof. Scholes,}

We submit the revised version of our manuscript ``Quantum-Coherent Spectral Engineering: Non-Markovian Dynamics in Photosynthetic Energy Transfer under Engineered Photon Baths'' (Manuscript ID: jz-2026-00994t) for consideration as a Letter in The Journal of Physical Chemistry Letters.

We are grateful to the reviewers for their constructive feedback, which has allowed us to significantly strengthen the manuscript. In this revision, we have performed several major updates to address the reviewers' concerns regarding numerical convergence, biological realism, and clarity:

\begin{enumerate}
\item \textbf{Rigorous convergence audit.} We have executed a new convergence audit up to hierarchy depth $L=10$ and $K=10$ Matsubara terms, providing a mathematical proof of convergence as requested by Reviewers 1 and 3.
\item \textbf{Physiological realism.} We have updated our simulation framework to incorporate the 12-mode spectral density for the FMO complex derived by Kleinekathöfer and Coker, addressing Reviewer 2's request for higher modeling fidelity.
\item \textbf{Visualization overhaul.} We have replotted all figures with a new 600 DPI publication theme, ensuring legible axes, units, and clear comparisons as requested by Reviewer 1 and 3.
\item \textbf{Theoretical reframing.} We have clarified the mechanism of "spectral bath engineering" using the polaron-frame formalism, distinguishing it from trivial state preparation.
\end{enumerate}

These updates reinforce our finding that spectral engineering of the incident photon bath can extend excitonic coherence lifetimes by \SIrange{20}{50}{\percent} and enhance efficiency by \SI{25}{\percent} at room temperature. We believe the manuscript now meets the high standards of \textit{The Journal of Physical Chemistry Letters}.

\textbf{Response to manuscript formatting request (28-Apr-2026).}
We have addressed all non-scientific formatting items as follows:
\begin{enumerate}
\item \textbf{Section headings removed.} The prohibited section headings (Introduction, Results and Discussion, Conclusion) have been removed. Major divisions are now indicated by bold run-in paragraph leads (e.g., \textbf{Theoretical Framework.}), and subsections use unnumbered \texttt{\textbackslash subsection*\{\}} headings, consistent with the JPCL Letter format.
\item \textbf{TOC Graphic included.} A new TOC Graphic is included via the \texttt{tocentry} environment in the \texttt{achemso} class, immediately preceding the document body, as required.
\item \textbf{Supporting Information.} The Supporting Information file has been checked and contains no highlighting, tracked changes, or colored text.
\end{enumerate}

\textbf{Cover art (Supplementary cover invitation, 28-Apr-2026).}
We are pleased to submit the TOC Graphic (\texttt{TOC\_Graphic.png}) for consideration as a supplementary journal cover. The image depicts the dual-band spectral filtering scheme applied to the FMO complex, with the engineered photon bath selectively populating vibronic resonances. We consent to both front and supplementary cover consideration.

All authors have approved this revised manuscript. We declare no conflicts of interest.

Thank you for your continued consideration of our work.

\closing{Sincerely,}

\fromsig{
Steve Cabrel Teguia Kouam\\
Department of Physics\\
University of Douala, Cameroon\\
Email: steve.teguia@facsciences-uy1.cm
}

\end{letter}
\end{document}
